% =============================================================================
% Abstract — reframed around the five empirical discoveries (F1--F5)
% Drop-in include: % =============================================================================
% Abstract — reframed around the five empirical discoveries (F1--F5)
% Drop-in include: % =============================================================================
% Abstract — reframed around the five empirical discoveries (F1--F5)
% Drop-in include: % =============================================================================
% Abstract — reframed around the five empirical discoveries (F1--F5)
% Drop-in include: \input{sections/abstract}
% Target: <= 250 words
% =============================================================================
\begin{abstract}
RL post-training of language models is now shaped by a handful of algorithms
(PPO, GRPO, DPO) and a growing fleet of frameworks (TRL, \tinker{}, OpenRLHF,
veRL), yet which choices actually drive performance remains empirically
unsettled.  We present \ours{}, a controlled study of 79 runs spanning 7 RL
libraries, 5 model families, and 0.6B--$\sim$1T parameters, designed to isolate
cause from confound on GSM8K, MATH-500, HumanEval, and tool-use.

Five discoveries reframe what should be measured and reported.
\textbf{(F1) ZVF as an early-warning diagnostic.}  Zero-Variance Fraction---the
share of GRPO groups with no within-group reward variance---correlates with
final reward (Pearson $r{=}{-}0.769$, $p{=}0.0008$).  A lagged regression
controlling for current reward shows ZVF's zero-order correlation is largely
explained by reward level (partial $r{=}0.059$ at lag 5, $p{=}0.096$), but
ZVF retains marginal independent signal at longer horizons (lag 10:
partial $r{=}0.080$, $p{=}0.045$) and is a practical triage tool: all runs
with sustained ZVF $\geq 80\%$ and reward $\leq 5\%$ in the first 5 steps
collapsed.
\textbf{(F2) Inverted-U in group size.}  $G{=}32$ maximizes gradient
utilization (54.5\%) and latest saturation onset (step 29); smaller and larger
$G$ are both worse, contradicting ``more rollouts is better.''
\textbf{(F3) Base vs.\ instruct dichotomy.}  Instruction tuning delivers a
$+0.922$ GSM8K delta (0.047$\to$0.969); Qwen3-30B-A3B-Instruct hits 100\%
last-10 where its base twin stalls at 32.5\%.  SFT, not RL, is the
trainability gate.
\textbf{(F4) PPO vs.\ GRPO may be model-dependent.}  In two single-seed
comparisons on matched compute, GRPO outperforms PPO on Qwen3-8B
(34.4\% vs.\ 22.5\% last-10) while PPO outperforms GRPO on Llama-3.1-8B
(97.5\% vs.\ 84.4\%).  We report this as a \emph{hypothesis} rather than a
conclusion: with only two model families and single-seed runs (MDE
$d{=}2.024$ at 80\% power), the interaction requires replication.
\textbf{(F5) Ceiling stability varies across architectures.}  At the 30-step
horizon, Nemotron-120B peaks 87.5\% then collapses to 16.2\% last-10, while
Kimi-K2 holds 100\%/80\% and Qwen3-235B-A22B holds 100\%/100\%---ceiling
stability is architecture- and pre-training-conditional, not a monotone
function of scale.  (All single-seed, early-training observations; asymptotic
convergence is not assessed.)
\textbf{(F6) Cold-Start Capability Threshold (validated).}  Critic-free
GRPO is blind whenever the probability of a mixed-reward group
$P(\text{usable})=1-(1-p_0)^{G}-p_0^{G}$ drops below the ZVF noise floor;
the familiar $p_0\!>\!1/G$ is the small-$p_0$ limit.  We validate this
predictor directly against every major task in the corpus
(strict vs.\ SFT-warmed tool-use, GSM8K, MATH-500, and binary vs.\
dense-reward HumanEval on the same Qwen3-8B base), recovering the
observed task hierarchy as a consequence of $P(\text{usable})$ rather
than of task semantics.
\textbf{(F7) Framework is a first-class experimental variable.}  On matched
model, task, and compute, the PPO-vs-GRPO winner flips between Qwen3-8B and
Llama-3.1-8B-Instruct, and TRL-vs-Tinker last-10 reward differs even under
the same nominal algorithm---library choice must be reported and ablated, not
treated as implementation detail.

Code, 79 run logs, and checkpoints:
\url{https://github.com/arvindcr4/tinker-rl-lab}.
\end{abstract}

% Target: <= 250 words
% =============================================================================
\begin{abstract}
RL post-training of language models is now shaped by a handful of algorithms
(PPO, GRPO, DPO) and a growing fleet of frameworks (TRL, \tinker{}, OpenRLHF,
veRL), yet which choices actually drive performance remains empirically
unsettled.  We present \ours{}, a controlled study of 79 runs spanning 7 RL
libraries, 5 model families, and 0.6B--$\sim$1T parameters, designed to isolate
cause from confound on GSM8K, MATH-500, HumanEval, and tool-use.

Five discoveries reframe what should be measured and reported.
\textbf{(F1) ZVF as an early-warning diagnostic.}  Zero-Variance Fraction---the
share of GRPO groups with no within-group reward variance---correlates with
final reward (Pearson $r{=}{-}0.769$, $p{=}0.0008$).  A lagged regression
controlling for current reward shows ZVF's zero-order correlation is largely
explained by reward level (partial $r{=}0.059$ at lag 5, $p{=}0.096$), but
ZVF retains marginal independent signal at longer horizons (lag 10:
partial $r{=}0.080$, $p{=}0.045$) and is a practical triage tool: all runs
with sustained ZVF $\geq 80\%$ and reward $\leq 5\%$ in the first 5 steps
collapsed.
\textbf{(F2) Inverted-U in group size.}  $G{=}32$ maximizes gradient
utilization (54.5\%) and latest saturation onset (step 29); smaller and larger
$G$ are both worse, contradicting ``more rollouts is better.''
\textbf{(F3) Base vs.\ instruct dichotomy.}  Instruction tuning delivers a
$+0.922$ GSM8K delta (0.047$\to$0.969); Qwen3-30B-A3B-Instruct hits 100\%
last-10 where its base twin stalls at 32.5\%.  SFT, not RL, is the
trainability gate.
\textbf{(F4) PPO vs.\ GRPO may be model-dependent.}  In two single-seed
comparisons on matched compute, GRPO outperforms PPO on Qwen3-8B
(34.4\% vs.\ 22.5\% last-10) while PPO outperforms GRPO on Llama-3.1-8B
(97.5\% vs.\ 84.4\%).  We report this as a \emph{hypothesis} rather than a
conclusion: with only two model families and single-seed runs (MDE
$d{=}2.024$ at 80\% power), the interaction requires replication.
\textbf{(F5) Ceiling stability varies across architectures.}  At the 30-step
horizon, Nemotron-120B peaks 87.5\% then collapses to 16.2\% last-10, while
Kimi-K2 holds 100\%/80\% and Qwen3-235B-A22B holds 100\%/100\%---ceiling
stability is architecture- and pre-training-conditional, not a monotone
function of scale.  (All single-seed, early-training observations; asymptotic
convergence is not assessed.)
\textbf{(F6) Cold-Start Capability Threshold (validated).}  Critic-free
GRPO is blind whenever the probability of a mixed-reward group
$P(\text{usable})=1-(1-p_0)^{G}-p_0^{G}$ drops below the ZVF noise floor;
the familiar $p_0\!>\!1/G$ is the small-$p_0$ limit.  We validate this
predictor directly against every major task in the corpus
(strict vs.\ SFT-warmed tool-use, GSM8K, MATH-500, and binary vs.\
dense-reward HumanEval on the same Qwen3-8B base), recovering the
observed task hierarchy as a consequence of $P(\text{usable})$ rather
than of task semantics.
\textbf{(F7) Framework is a first-class experimental variable.}  On matched
model, task, and compute, the PPO-vs-GRPO winner flips between Qwen3-8B and
Llama-3.1-8B-Instruct, and TRL-vs-Tinker last-10 reward differs even under
the same nominal algorithm---library choice must be reported and ablated, not
treated as implementation detail.

Code, 79 run logs, and checkpoints:
\url{https://github.com/arvindcr4/tinker-rl-lab}.
\end{abstract}

% Target: <= 250 words
% =============================================================================
\begin{abstract}
RL post-training of language models is now shaped by a handful of algorithms
(PPO, GRPO, DPO) and a growing fleet of frameworks (TRL, \tinker{}, OpenRLHF,
veRL), yet which choices actually drive performance remains empirically
unsettled.  We present \ours{}, a controlled study of 79 runs spanning 7 RL
libraries, 5 model families, and 0.6B--$\sim$1T parameters, designed to isolate
cause from confound on GSM8K, MATH-500, HumanEval, and tool-use.

Five discoveries reframe what should be measured and reported.
\textbf{(F1) ZVF as an early-warning diagnostic.}  Zero-Variance Fraction---the
share of GRPO groups with no within-group reward variance---correlates with
final reward (Pearson $r{=}{-}0.769$, $p{=}0.0008$).  A lagged regression
controlling for current reward shows ZVF's zero-order correlation is largely
explained by reward level (partial $r{=}0.059$ at lag 5, $p{=}0.096$), but
ZVF retains marginal independent signal at longer horizons (lag 10:
partial $r{=}0.080$, $p{=}0.045$) and is a practical triage tool: all runs
with sustained ZVF $\geq 80\%$ and reward $\leq 5\%$ in the first 5 steps
collapsed.
\textbf{(F2) Inverted-U in group size.}  $G{=}32$ maximizes gradient
utilization (54.5\%) and latest saturation onset (step 29); smaller and larger
$G$ are both worse, contradicting ``more rollouts is better.''
\textbf{(F3) Base vs.\ instruct dichotomy.}  Instruction tuning delivers a
$+0.922$ GSM8K delta (0.047$\to$0.969); Qwen3-30B-A3B-Instruct hits 100\%
last-10 where its base twin stalls at 32.5\%.  SFT, not RL, is the
trainability gate.
\textbf{(F4) PPO vs.\ GRPO may be model-dependent.}  In two single-seed
comparisons on matched compute, GRPO outperforms PPO on Qwen3-8B
(34.4\% vs.\ 22.5\% last-10) while PPO outperforms GRPO on Llama-3.1-8B
(97.5\% vs.\ 84.4\%).  We report this as a \emph{hypothesis} rather than a
conclusion: with only two model families and single-seed runs (MDE
$d{=}2.024$ at 80\% power), the interaction requires replication.
\textbf{(F5) Ceiling stability varies across architectures.}  At the 30-step
horizon, Nemotron-120B peaks 87.5\% then collapses to 16.2\% last-10, while
Kimi-K2 holds 100\%/80\% and Qwen3-235B-A22B holds 100\%/100\%---ceiling
stability is architecture- and pre-training-conditional, not a monotone
function of scale.  (All single-seed, early-training observations; asymptotic
convergence is not assessed.)
\textbf{(F6) Cold-Start Capability Threshold (validated).}  Critic-free
GRPO is blind whenever the probability of a mixed-reward group
$P(\text{usable})=1-(1-p_0)^{G}-p_0^{G}$ drops below the ZVF noise floor;
the familiar $p_0\!>\!1/G$ is the small-$p_0$ limit.  We validate this
predictor directly against every major task in the corpus
(strict vs.\ SFT-warmed tool-use, GSM8K, MATH-500, and binary vs.\
dense-reward HumanEval on the same Qwen3-8B base), recovering the
observed task hierarchy as a consequence of $P(\text{usable})$ rather
than of task semantics.
\textbf{(F7) Framework is a first-class experimental variable.}  On matched
model, task, and compute, the PPO-vs-GRPO winner flips between Qwen3-8B and
Llama-3.1-8B-Instruct, and TRL-vs-Tinker last-10 reward differs even under
the same nominal algorithm---library choice must be reported and ablated, not
treated as implementation detail.

Code, 79 run logs, and checkpoints:
\url{https://github.com/arvindcr4/tinker-rl-lab}.
\end{abstract}

% Target: <= 250 words
% =============================================================================
\begin{abstract}
RL post-training of language models is now shaped by a handful of algorithms
(PPO, GRPO, DPO) and a growing fleet of frameworks (TRL, \tinker{}, OpenRLHF,
veRL), yet which choices actually drive performance remains empirically
unsettled.  We present \ours{}, a controlled study of 79 runs spanning 7 RL
libraries, 5 model families, and 0.6B--$\sim$1T parameters, designed to isolate
cause from confound on GSM8K, MATH-500, HumanEval, and tool-use.

Five discoveries reframe what should be measured and reported.
\textbf{(F1) ZVF as an early-warning diagnostic.}  Zero-Variance Fraction---the
share of GRPO groups with no within-group reward variance---correlates with
final reward (Pearson $r{=}{-}0.769$, $p{=}0.0008$).  A lagged regression
controlling for current reward shows ZVF's zero-order correlation is largely
explained by reward level (partial $r{=}0.059$ at lag 5, $p{=}0.096$), but
ZVF retains marginal independent signal at longer horizons (lag 10:
partial $r{=}0.080$, $p{=}0.045$) and is a practical triage tool: all runs
with sustained ZVF $\geq 80\%$ and reward $\leq 5\%$ in the first 5 steps
collapsed.
\textbf{(F2) Inverted-U in group size.}  $G{=}32$ maximizes gradient
utilization (54.5\%) and latest saturation onset (step 29); smaller and larger
$G$ are both worse, contradicting ``more rollouts is better.''
\textbf{(F3) Base vs.\ instruct dichotomy.}  Instruction tuning delivers a
$+0.922$ GSM8K delta (0.047$\to$0.969); Qwen3-30B-A3B-Instruct hits 100\%
last-10 where its base twin stalls at 32.5\%.  SFT, not RL, is the
trainability gate.
\textbf{(F4) PPO vs.\ GRPO may be model-dependent.}  In two single-seed
comparisons on matched compute, GRPO outperforms PPO on Qwen3-8B
(34.4\% vs.\ 22.5\% last-10) while PPO outperforms GRPO on Llama-3.1-8B
(97.5\% vs.\ 84.4\%).  We report this as a \emph{hypothesis} rather than a
conclusion: with only two model families and single-seed runs (MDE
$d{=}2.024$ at 80\% power), the interaction requires replication.
\textbf{(F5) Ceiling stability varies across architectures.}  At the 30-step
horizon, Nemotron-120B peaks 87.5\% then collapses to 16.2\% last-10, while
Kimi-K2 holds 100\%/80\% and Qwen3-235B-A22B holds 100\%/100\%---ceiling
stability is architecture- and pre-training-conditional, not a monotone
function of scale.  (All single-seed, early-training observations; asymptotic
convergence is not assessed.)
\textbf{(F6) Cold-Start Capability Threshold (validated).}  Critic-free
GRPO is blind whenever the probability of a mixed-reward group
$P(\text{usable})=1-(1-p_0)^{G}-p_0^{G}$ drops below the ZVF noise floor;
the familiar $p_0\!>\!1/G$ is the small-$p_0$ limit.  We validate this
predictor directly against every major task in the corpus
(strict vs.\ SFT-warmed tool-use, GSM8K, MATH-500, and binary vs.\
dense-reward HumanEval on the same Qwen3-8B base), recovering the
observed task hierarchy as a consequence of $P(\text{usable})$ rather
than of task semantics.
\textbf{(F7) Framework is a first-class experimental variable.}  On matched
model, task, and compute, the PPO-vs-GRPO winner flips between Qwen3-8B and
Llama-3.1-8B-Instruct, and TRL-vs-Tinker last-10 reward differs even under
the same nominal algorithm---library choice must be reported and ablated, not
treated as implementation detail.

Code, 79 run logs, and checkpoints:
\url{https://github.com/arvindcr4/tinker-rl-lab}.
\end{abstract}
